% ============================================================================
% 公共样式与跨平台排版配置（参考 NeurIPS 论文排版比例，纯黑正文字体 + 彩色配图）
% ============================================================================

% --- 页面尺寸与页边距（NeurIPS 单栏学术论文版芯比例）---
\usepackage[a4paper, top=2.4cm, bottom=2.4cm, left=2.7cm, right=2.7cm, headheight=14pt]{geometry}

% --- 跨平台中英文字体配置（自动适配 Windows / macOS / Linux & Overleaf）---
\usepackage{fontspec}
\IfFontExistsTF{Times New Roman}{
  \setmainfont{Times New Roman}
}{
  \IfFontExistsTF{TeX Gyre Termes}{
    \setmainfont{TeX Gyre Termes}
  }{}
}
\IfFontExistsTF{Arial}{
  \setsansfont{Arial}
}{
  \IfFontExistsTF{TeX Gyre Heros}{
    \setsansfont{TeX Gyre Heros}
  }{}
}
\IfFontExistsTF{Menlo}{
  \setmonofont[Scale=0.86]{Menlo}
}{
  \IfFontExistsTF{DejaVu Sans Mono}{
    \setmonofont[Scale=0.85]{DejaVu Sans Mono}
  }{}
}

% 正文中文字体：宋体（优先 SimSun -> Songti SC -> STSong -> FandolSong）
\IfFontExistsTF{SimSun}{
  \setCJKmainfont[AutoFakeBold=2.2, ItalicFont=KaiTi]{SimSun}
}{
  \IfFontExistsTF{Songti SC}{
    \setCJKmainfont[BoldFont=Songti SC Bold, ItalicFont=Kaiti SC]{Songti SC}
  }{
    \IfFontExistsTF{STSong}{
      \setCJKmainfont[AutoFakeBold=2.2, ItalicFont=Kaiti SC]{STSong}
    }{
      \setCJKmainfont[AutoFakeBold=2.2]{FandolSong}
    }
  }
}

% 无衬线中文字体：黑体
\IfFontExistsTF{SimHei}{
  \setCJKsansfont[AutoFakeBold=1.8]{SimHei}
}{
  \IfFontExistsTF{Heiti SC}{
    \setCJKsansfont[BoldFont=Heiti SC Medium]{Heiti SC}
  }{
    \IfFontExistsTF{STHeiti}{
      \setCJKsansfont[AutoFakeBold=1.8]{STHeiti}
    }{
      \setCJKsansfont[AutoFakeBold=1.8]{FandolHei}
    }
  }
}

% 等宽中文字体与楷体定义
\IfFontExistsTF{FangSong}{
  \setCJKmonofont{FangSong}
}{
  \IfFontExistsTF{Kaiti SC}{
    \setCJKmonofont{Kaiti SC}
  }{
    \setCJKmonofont{FandolFang}
  }
}
\IfFontExistsTF{KaiTi}{
  \newCJKfontfamily\kaiti{KaiTi}
}{
  \IfFontExistsTF{Kaiti SC}{
    \newCJKfontfamily\kaiti{Kaiti SC}
  }{
    \newCJKfontfamily\kaiti{FandolKai}
  }
}

% --- 数学、表格、图形与算法宏包 ---
\usepackage{amsmath, amssymb, amsthm}
\usepackage{booktabs}
\usepackage{tabularx}
\usepackage{multirow}
\usepackage{array}
\usepackage{graphicx}
\usepackage{xcolor}
\usepackage{float}
\usepackage{caption}
\usepackage{subcaption}
\usepackage{enumitem}
\usepackage{setspace}
\usepackage{tikz}
\usetikzlibrary{arrows.meta, positioning, shapes.geometric, fit, backgrounds, calc}
\usepackage[most]{tcolorbox}

% --- 参考文献与超链接 ---
\usepackage[numbers, sort&compress]{natbib}
\usepackage[hidelinks, bookmarksnumbered=true]{hyperref}

% --- 配图专用色彩定义（仅用于 TikZ 图形内部，标题与正文全部为纯黑）---
\definecolor{PrimaryNavy}{HTML}{1B365D}
\definecolor{AccentBlue}{HTML}{2B6CB0}
\definecolor{SoftBlueBg}{HTML}{EDF2F7}
\definecolor{AccentTeal}{HTML}{2C7A7B}
\definecolor{SoftTealBg}{HTML}{E6FFFA}
\definecolor{WarmAmber}{HTML}{B7791F}
\definecolor{SoftAmberBg}{HTML}{FEFCBF}
\definecolor{BorderGray}{HTML}{CBD5E0}
\definecolor{MutedText}{HTML}{4A5568}

% --- 紧凑学术行距与段落间距（对齐 NeurIPS 10pt 论文视觉密度）---
\linespread{1.18}
\setlength{\parskip}{3.5pt}
\setlength{\emergencystretch}{2em}
\setlist[itemize]{itemsep=1.5pt, topsep=3pt, parsep=0pt}
\setlist[enumerate]{itemsep=1.5pt, topsep=3pt, parsep=0pt}

% --- 章节标题格式（参考 NeurIPS 分级字号：纯黑加粗，紧凑克制）---
% 正文基准为 10pt (五号 10.5pt)：
% 一级标题 \section: 12pt 加粗
% 二级标题 \subsection: 10.5pt 加粗
% 三级标题 \subsubsection: 10pt 加粗
\ctexset{
  section = {
    format = \fontsize{12pt}{15pt}\selectfont\bfseries,
    name = {,},
    number = \arabic{section},
    aftername = \hspace{0.55em},
    beforeskip = 12pt plus 3pt minus 2pt,
    afterskip = 5pt plus 1pt minus 1pt,
  },
  subsection = {
    format = \fontsize{10.5pt}{13.5pt}\selectfont\bfseries,
    aftername = \hspace{0.45em},
    beforeskip = 9pt plus 2pt minus 1pt,
    afterskip = 4pt plus 1pt minus 1pt,
  },
  subsubsection = {
    format = \fontsize{10pt}{12.5pt}\selectfont\bfseries,
    aftername = \hspace{0.4em},
    beforeskip = 6pt plus 1pt minus 1pt,
    afterskip = 3pt plus 1pt minus 1pt,
  }
}

% --- 图表标题格式（纯黑紧凑 9pt）---
\captionsetup{
  font={footnotesize},
  labelfont={bf},
  labelsep=colon,
  skip=5pt
}

% --- 页眉与页脚严格规则 ---
% 第 1 页：顶部居中显示《人工智能基础》大作业 + 单横线（对应 P2）
% 第 2 页及以后：完全取消页眉文字与页眉横线（删去 P4），仅保留底部居中页码
\usepackage{fancyhdr}

\fancypagestyle{firstpage}{
  \fancyhf{}
  \fancyhead[C]{\fontsize{9.5pt}{11pt}\selectfont 《人工智能基础》大作业}
  \renewcommand{\headrulewidth}{0.6pt}
  \renewcommand{\footrulewidth}{0pt}
  \fancyfoot[C]{\small\thepage}
}

\fancypagestyle{mainbody}{
  \fancyhf{}
  \renewcommand{\headrulewidth}{0pt}
  \renewcommand{\footrulewidth}{0pt}
  \fancyfoot[C]{\small\thepage}
}

\pagestyle{mainbody}

% --- 队友留白占位框（纯黑文字配浅灰虚线框，标明对应队友提交模块）---
\newtcolorbox{teammatebox}[1]{
  enhanced,
  breakable,
  colback=black!2,
  colframe=black!45,
  coltitle=black,
  fonttitle=\bfseries\footnotesize,
  title={队友负责模块留白（待队友补充）：#1},
  boxrule=0.5pt,
  arc=2pt,
  left=8pt,
  right=8pt,
  top=5pt,
  bottom=5pt,
  fontupper=\footnotesize,
  attach boxed title to top left={xshift=8pt, yshift=-2mm},
  boxed title style={colback=white, colframe=black!45, boxrule=0.5pt, arc=2pt}
}

% --- 论文标题与双人居中署名区（参考上传 NeurIPS 论文标题与作者字号比例，纯黑）---
% 标题：上下双横线（上粗 2.5pt、下细 0.8pt），字号 16pt 加粗居中
% 学号与姓名：与 NeurIPS 作者栏一致采用 10pt 加粗居中，两行紧凑排列
\newcommand{\MakePaperTitle}[5]{%
  \thispagestyle{firstpage}%
  \begin{center}
    \vspace*{2pt}
    \noindent\rule{\textwidth}{2.4pt}\par
    \vspace{6pt}
    {\fontsize{15.5pt}{21pt}\selectfont\bfseries #1 \par}
    \vspace{4pt}
    \noindent\rule{\textwidth}{0.7pt}\par
    \vspace{9pt}
    {\fontsize{10pt}{14.5pt}\selectfont\bfseries
      学号：#2，姓名：#3 \\[2pt]
      学号：#4，姓名：#5 \par
    }
    \vspace{8pt}
  \end{center}
}

% --- 代码标识符行内等宽排版 ---
\newcommand{\code}[1]{\texttt{#1}}
